\section{Introduction}
\label{sec:intro}

The game of 20 Questions distills a fundamental challenge in information seeking: given a set of possibilities, ask questions that narrow the space as efficiently as possible.
The optimal strategy is binary search---each question should split the remaining candidates into two equal halves, maximizing information gain at every step.
This principle underlies not only parlor games but also medical diagnosis, fault isolation, and interactive AI systems that must identify user intent through clarifying questions.

Given the rapid advances in large language model (LLM) capabilities, a natural question arises: \textbf{can LLMs generate yes/no questions that split a set of items in half?}
If so, they could serve as efficient question generators in interactive systems without any task-specific training.
If not, understanding \emph{when} and \emph{why} they fail would inform the design of better prompting strategies and training objectives.

Prior work has evaluated LLMs as players in 20 Questions games, finding that zero-shot models achieve expected information gain (EIG) of only 0.29--0.35 \citep{mazzaccara2024learning}, far from the optimal 1.0 bit.
However, these studies tested older models (GPT-3.5, LLaMA~2) and embedded question quality within multi-turn game performance, making it difficult to isolate the core capability.
No prior work has systematically evaluated whether modern LLMs (GPT-4.1, released March 2025) can generate a \emph{single} question that achieves maximum entropy over a given item set, or what properties of the set make this task difficult.

We address this gap with the first systematic evaluation of single-question split quality across modern LLMs, multiple prompt strategies, and diverse set compositions.
Our key findings include:

\begin{itemize}[leftmargin=*,itemsep=0pt,topsep=0pt]
    \item We show that the capability for optimal splitting \emph{exists} but is not default behavior: \gptfull with chain-of-thought prompting achieves binary entropy of 1.000 on structured sets, compared to 0.027 with basic prompting on homogeneous sets---a gap of 0.973.
    \item We identify set composition as the primary difficulty factor, establishing a clear difficulty ranking from mixed-category sets (100\% perfect splits) to homogeneous sets (65\% perfect splits).
    \item We demonstrate that \gptmini produces more naturally discriminative questions than \gptfull under basic prompting, revealing that stronger models do not automatically ask better questions.
    \item We evaluate three prompt strategies across 8 datasets and 1,049 trials, finding that chain-of-thought prompting is the most reliable approach for eliciting optimal splits.
\end{itemize}

The rest of the paper is organized as follows. \Secref{sec:related} reviews related work on LLM question generation and information-theoretic evaluation. \Secref{sec:method} describes our methodology, datasets, and metrics. \Secref{sec:results} presents experimental results and analysis. \Secref{sec:discussion} discusses implications and limitations, and \secref{sec:conclusion} concludes.
